\documentclass[aspectratio=169]{beamer}
\usetheme{Madrid}
\usecolortheme{default}
\usepackage{booktabs}
\usepackage{graphicx}
\usepackage{listings}
\usepackage{xcolor}
\usepackage{tikz}

% Custom colours
\definecolor{ntnu}{RGB}{0, 80, 158}
\definecolor{lightgray}{RGB}{240, 240, 240}
\definecolor{darkgreen}{RGB}{0, 130, 0}

\setbeamercolor{frametitle}{bg=ntnu, fg=white}
\setbeamercolor{title}{bg=ntnu, fg=white}
\setbeamercolor{block title}{bg=ntnu, fg=white}
\setbeamercolor{structure}{fg=ntnu}

\lstset{
  basicstyle=\ttfamily\footnotesize,
  backgroundcolor=\color{lightgray},
  frame=single,
  breaklines=true,
  keywordstyle=\color{ntnu}\bfseries,
  commentstyle=\color{gray},
  language=Python
}

% Title info
\title{TDT4265 Mini-Project}
\subtitle{Snow Pole Detection for Autonomous Driving}
\author{Gustav}
\institute{NTNU -- Computer Vision and Deep Learning}
\date{April 2026}

% -------------------------------------------------------
\begin{document}

\begin{frame}
  \titlepage
\end{frame}

\begin{frame}{Outline}
  \tableofcontents
\end{frame}

% =======================================================
\section{Methods \& Models}
% =======================================================

\begin{frame}{Overview -- Two Architectures}
  \begin{columns}
    \begin{column}{0.48\textwidth}
      \begin{block}{Model 1 -- YOLOv8s}
        \begin{itemize}
          \item Single-stage detector
          \item You Only Look Once -- one forward pass
          \item Designed for real-time edge deployment
          \item \textbf{Dataset:} Road poles iPhone\\(942 train / 261 val)
        \end{itemize}
      \end{block}
    \end{column}
    \begin{column}{0.48\textwidth}
      \begin{block}{Model 2 -- Faster R-CNN}
        \begin{itemize}
          \item Two-stage detector
          \item Region Proposal Network + classification
          \item ResNet-50 + FPN backbone
          \item \textbf{Dataset:} roadpoles v1\\(322 train / 92 val)
        \end{itemize}
      \end{block}
    \end{column}
  \end{columns}
  \vspace{0.5em}
  \centering
  \textit{Both models fine-tuned from ImageNet/COCO pretrained weights}
\end{frame}

% -------------------------------------------------------
\begin{frame}{Model 1 -- YOLOv8s Architecture}
  \begin{columns}
    \begin{column}{0.55\textwidth}
      \textbf{Why YOLOv8s?}
      \begin{itemize}
        \item Poles occupy $\approx$2\% of image width -- need enough model capacity
        \item YOLOv8\textbf{n} (nano) too weak for small objects
        \item YOLOv8\textbf{s} (small, 11.2M params) -- best speed/accuracy tradeoff
        \item Must run on edge device in real-time
      \end{itemize}
      \vspace{0.5em}
      \textbf{Key training decisions:}
      \begin{itemize}
        \item \texttt{imgsz=1280} -- more pixels per pole
        \item 150 epochs with early stopping (patience=30)
        \item Mosaic, scale, colour jitter augmentations
      \end{itemize}
    \end{column}
    \begin{column}{0.42\textwidth}
      \begin{block}{Model stats}
        \begin{itemize}
          \item Parameters: 11.2M
          \item Input size: $1280\times1280$
          \item Classes: 1 (pole)
          \item Pretrained: COCO
          \item Batch size: 8
          \item Optimiser: auto (AdamW)
        \end{itemize}
      \end{block}
    \end{column}
  \end{columns}
\end{frame}

% -------------------------------------------------------
\begin{frame}[fragile]{Model 1 -- YOLOv8s Training Config}
  \begin{block}{Training code}
\begin{lstlisting}[language=Python]
model = YOLO("yolov8s.pt")
model.train(
    data="data.yaml",
    epochs=150,
    imgsz=1280,
    batch=8,
    patience=30,
    hsv_h=0.015, hsv_s=0.7, hsv_v=0.4,
    fliplr=0.5, mosaic=1.0, scale=0.5,
    translate=0.1,
    device=0
)
\end{lstlisting}
  \end{block}
  \begin{itemize}
    \item Batch size reduced to 8 due to \texttt{imgsz=1280} VRAM requirements
    \item Fine-tuned from COCO pretrained weights
    \item Trained on NVIDIA RTX 4090 (Cybele lab)
  \end{itemize}
\end{frame}

% -------------------------------------------------------
\begin{frame}{Model 2 -- Faster R-CNN Architecture}
  \begin{columns}
    \begin{column}{0.55\textwidth}
      \textbf{Two-stage detection pipeline:}
      \begin{enumerate}
        \item \textbf{Stage 1 -- RPN:} Scans feature map, proposes candidate regions containing objects
        \item \textbf{Stage 2 -- ROI Head:} Classifies each proposal and refines the bounding box
      \end{enumerate}
      \vspace{0.5em}
      \textbf{FPN backbone:} Extracts features at multiple scales simultaneously -- important for poles at different distances
      \vspace{0.5em}
      \textbf{Key difference from YOLO:} Two passes = slower but theoretically more accurate on small objects
    \end{column}
    \begin{column}{0.42\textwidth}
      \begin{block}{Architecture}
        \begin{itemize}
          \item Backbone: ResNet-50 + FPN
          \item Pretrained: COCO
          \item Head: replaced for 2 classes (background + pole)
          \item Optimiser: SGD, lr=0.005
          \item Scheduler: StepLR
          \item Epochs: 100
        \end{itemize}
      \end{block}
    \end{column}
  \end{columns}
\end{frame}

% -------------------------------------------------------
\begin{frame}[fragile]{Model 2 -- Faster R-CNN Training Config}
  \begin{block}{Custom dataset class -- YOLO labels to Faster R-CNN format}
\begin{lstlisting}[language=Python]
# Convert normalised YOLO format to absolute pixel coords
x1 = (cx - bw/2) * W
y1 = (cy - bh/2) * H
x2 = (cx + bw/2) * W
y2 = (cy + bh/2) * H
# Labels: 0 = background, 1 = pole
\end{lstlisting}
  \end{block}
  \begin{block}{Augmentations (train only)}
\begin{lstlisting}[language=Python]
train_transform = T.Compose([
    T.ToTensor(),
    T.ColorJitter(brightness=0.4, contrast=0.4,
                  saturation=0.3, hue=0.1),
    T.RandomHorizontalFlip(p=0.5),
])
\end{lstlisting}
  \end{block}
\end{frame}

% =======================================================
\section{Results}
% =======================================================

\begin{frame}{Results -- Validation Set Metrics}
  \begin{table}
    \centering
    \renewcommand{\arraystretch}{1.4}
    \begin{tabular}{lcccc}
      \toprule
      \textbf{Model} & \textbf{Precision} & \textbf{Recall} & \textbf{mAP@50} & \textbf{mAP@0.5:0.95} \\
      \midrule
      YOLOv8s        & 0.957 & 0.976 & \textbf{0.990} & \textbf{0.788} \\
      Faster R-CNN   & 0.626 & 0.679 & 0.993 & 0.626 \\
      \bottomrule
    \end{tabular}
    \caption{Validation set results on respective datasets}
  \end{table}
  \vspace{0.5em}
  \begin{itemize}
    \item Both models achieve \textbf{mAP@50 $>$ 0.99} -- poles are reliably detected
    \item YOLO significantly better on \textbf{mAP@0.5:0.95} (0.788 vs 0.626)
    \item Faster R-CNN higher Precision but lower Recall
  \end{itemize}
\end{frame}

% -------------------------------------------------------
\begin{frame}{Results -- Leaderboard}
  \begin{columns}
    \begin{column}{0.48\textwidth}
      \begin{block}{Road poles iPhone leaderboard}
        \begin{itemize}
          \item Submitted YOLOv8s predictions
          \item Score: \textbf{$\approx$69\%}
        \end{itemize}
      \end{block}
    \end{column}
    \begin{column}{0.48\textwidth}
      \begin{block}{roadpoles v1 leaderboard}
        \begin{itemize}
          \item Submitted Faster R-CNN predictions
          \item Score: \textbf{$\approx$66\%}
        \end{itemize}
      \end{block}
    \end{column}
  \end{columns}
  \vspace{1em}
  \textbf{Note:} Test set scores lower than validation -- expected due to unseen weather/lighting conditions in test set
\end{frame}

% -------------------------------------------------------
\begin{frame}{Results -- mAP@50 vs mAP@0.5:0.95 Gap}
  \begin{columns}
    \begin{column}{0.55\textwidth}
      The gap between the two metrics tells us how \textbf{precisely the boxes fit} the poles:
      \vspace{0.5em}
      \begin{itemize}
        \item \textbf{mAP@50:} Box needs to overlap ground truth by $\geq$50\% -- both models pass easily
        \item \textbf{mAP@0.5:0.95:} Averaged over IoU thresholds 0.50 to 0.95 -- requires tight localisation
      \end{itemize}
      \vspace{0.5em}
      YOLO produces \textbf{tighter bounding boxes} -- likely due to larger training set and higher image resolution
    \end{column}
    \begin{column}{0.42\textwidth}
      \begin{table}
        \renewcommand{\arraystretch}{1.4}
        \begin{tabular}{lcc}
          \toprule
          & \textbf{@50} & \textbf{@0.5:0.95} \\
          \midrule
          YOLOv8s      & 0.990 & 0.788 \\
          Faster R-CNN & 0.993 & 0.626 \\
          \bottomrule
        \end{tabular}
      \end{table}
      \vspace{0.5em}
      \textcolor{darkgreen}{Gap for YOLO: 0.202}\\
      \textcolor{red}{Gap for RCNN: 0.367}
    \end{column}
  \end{columns}
\end{frame}

% =======================================================
\section{Discussion}
% =======================================================

\begin{frame}{Discussion -- Why YOLO Outperformed Faster R-CNN}
  \begin{block}{Three key factors}
    \begin{enumerate}
      \item \textbf{Dataset size:} YOLO trained on 942 images vs 322 for Faster R-CNN. More data = better generalisation, especially for small objects.
      \item \textbf{Image resolution:} YOLO used \texttt{imgsz=1280} vs default 800 for Faster R-CNN. More pixels per pole directly improves detection.
      \item \textbf{Augmentation:} YOLO used mosaic augmentation which creates composite training images -- effectively increases dataset diversity.
    \end{enumerate}
  \end{block}
  \vspace{0.3em}
  \textit{Faster R-CNN is theoretically stronger on small objects due to FPN + two-stage design, but needs more data and tuning to demonstrate this advantage.}
\end{frame}

% -------------------------------------------------------
\begin{frame}{Discussion -- Limitations \& Future Work}
  \begin{columns}
    \begin{column}{0.48\textwidth}
      \textbf{Current limitations:}
      \begin{itemize}
        \item Faster R-CNN trained on smaller dataset -- unfair comparison
        \item Test set gap suggests domain shift (weather, lighting)
        \item Faster R-CNN too slow for real-time edge deployment as-is
        \item No pseudo-labelling of unlabelled RoadPoles-MSJ data
      \end{itemize}
    \end{column}
    \begin{column}{0.48\textwidth}
      \textbf{Future improvements:}
      \begin{itemize}
        \item Increase Faster R-CNN \texttt{imgsz} to 1024/1280
        \item Pseudo-label RoadPoles-MSJ with best YOLO model
        \item Try MobileNet backbone for faster inference
        \item Test time augmentation (TTA) for better test scores
        \item Ensemble YOLO + Faster R-CNN predictions
      \end{itemize}
    \end{column}
  \end{columns}
\end{frame}

% =======================================================
\section{Key Learning Points}
% =======================================================

\begin{frame}{Key Learning Points}
  \begin{block}{1 -- Resolution matters more than model size}
    Increasing \texttt{imgsz} from 640 to 1280 had a larger impact than switching from YOLOv8n to YOLOv8s. For tiny objects like snow poles, giving the model more pixels to work with is the single most effective improvement.
  \end{block}
  \begin{block}{2 -- Understanding two-stage detection}
    Implementing Faster R-CNN manually revealed how the RPN and ROI pooling work together -- a much deeper understanding than using YOLO through a high-level API. The label format conversion (YOLO normalised to absolute pixel coordinates) highlighted how different frameworks expect data.
  \end{block}
  \begin{block}{3 -- Applied deep learning is mostly engineering}
    A significant portion of the project was path issues, yaml configuration, CUDA driver mismatches, and dataset format debugging. These are the real-world challenges in any CV deployment scenario.
  \end{block}
\end{frame}

% =======================================================
\section{Sustainability}
% =======================================================

\begin{frame}{Sustainability -- Compute Energy}
  \begin{columns}
    \begin{column}{0.55\textwidth}
      \textbf{Hardware used:}
      \begin{itemize}
        \item NVIDIA RTX 4090 (Cybele lab)
        \item TDP: 450W
      \end{itemize}
      \vspace{0.5em}
      \textbf{Approximate training time:}
      \begin{itemize}
        \item YOLOv8s (150 epochs, imgsz=1280): $\approx$2.5h
        \item Faster R-CNN (100 epochs): $\approx$1.5h
        \item Total: $\approx$\textbf{4 hours}
      \end{itemize}
      \vspace{0.5em}
      \textbf{Energy consumed:}
      \[E = 0.45\,\text{kW} \times 4\,\text{h} = \mathbf{1.8\,\text{kWh}}\]
    \end{column}
    \begin{column}{0.42\textwidth}
      \begin{block}{How long could we code?}
        \vspace{0.3em}
        MacBook Pro M3 consumption:\\
        \textbf{$\approx$45 W under load}
        \vspace{0.5em}
        \[t = \frac{1800\,\text{Wh}}{45\,\text{W}} = \mathbf{40\,\text{hours}}\]
        \vspace{0.3em}
        \centering
        $\approx$ a full work week of laptop use
        \vspace{0.3em}
        \textit{Relatively low footprint -- fine-tuning from pretrained weights avoided training from scratch}
      \end{block}
    \end{column}
  \end{columns}
\end{frame}

% -------------------------------------------------------
\begin{frame}{Summary}
  \begin{table}
    \centering
    \renewcommand{\arraystretch}{1.4}
    \begin{tabular}{lcccc}
      \toprule
      \textbf{Model} & \textbf{mAP@50} & \textbf{mAP@0.5:0.95} & \textbf{Speed} & \textbf{Edge-ready} \\
      \midrule
      YOLOv8s      & 0.990 & 0.788 & Fast  & \textcolor{darkgreen}{\textbf{Yes}} \\
      Faster R-CNN & 0.993 & 0.626 & Slow  & \textcolor{red}{\textbf{No}}  \\
      \bottomrule
    \end{tabular}
  \end{table}
  \vspace{1em}
  \begin{itemize}
    \item \textbf{Best model for deployment:} YOLOv8s -- higher mAP@0.5:0.95, real-time capable
    \item \textbf{Biggest improvement lever:} image resolution for small object detection
    \item \textbf{Energy cost:} 1.8 kWh total $\approx$ 12 km in a Tesla Model Y
  \end{itemize}
  \vspace{0.5em}
  \centering
  \textit{Thank you!}
\end{frame}

\end{document}